\uuid{ZjQr}
\titre{Module/argument}
\auteur{bodin ; hanani ; mzari}
\organisation{exo7}
\licence{CC BY-NC-SA 4.0}
\datecreate{2020-01-20}
\niveau{L1}
\module{Algèbre}
\chapitre{Nombres complexes}
\sousChapitre{Forme cartésienne, forme polaire}
\difficulte{3}
\theme{Écritures algébrique et géométrique, module, argument}
\exo7qcmid{lille-1.nombres-complexes.012}

\contenu{
\begin{question}
\qclassification{104.01}
\qtags{motcle=module/argument}
Par définition, si  $x,y \in \Rr, \, e^{x+iy} = e^x \cdot e^{iy}= e^x (\cos y +i \sin y)$.

Soit $z=e^{e^{i\theta}}$, où $\theta$ est un réel. Quelles sont les assertions vraies ?
\begin{answers}
    \bad{$|z|=1 $}

    \good{$|z|=e^{\cos \theta} $}

    \bad{$\arg z = \theta \,  [2\pi]$}

    \good{$\arg z = \sin \theta \, [2\pi]$}

\end{answers}
\begin{explanations}
$z= e^{\cos  \theta + i \sin \theta}= e^{\cos\theta}\cdot e^{i \sin \theta}.  $ Donc $|z|=e^{\cos \theta} $ et $\arg z = \sin \theta \, [2\pi]$.
\end{explanations}
\end{question}
}
